\documentclass[11pt,a4paper]{article}

\usepackage[T2A]{fontenc}
\usepackage[utf8]{inputenc}
\usepackage[russian]{babel}
\usepackage{geometry}
\usepackage{microtype}
\usepackage{listings}
\usepackage{xcolor}
\usepackage{hyperref}
\usepackage{amsmath}
\usepackage{amssymb}
\usepackage{graphicx}
\usepackage{float}
\usepackage{booktabs}
\usepackage{enumitem}
\usepackage{titlesec}
\usepackage{fancyhdr}
\usepackage{setspace}
\usepackage{parskip}
\usepackage{caption}
\usepackage{subcaption}
\usepackage{mdframed}

% Поля — для комфортного чтения на экране 13"
\geometry{
  a4paper,
  left=2.8cm,
  right=2.8cm,
  top=3cm,
  bottom=3cm
}

% Гиперссылки
\hypersetup{
  colorlinks=true,
  linkcolor=blue!60!black,
  citecolor=green!50!black,
  urlcolor=blue!70!black,
  pdftitle={Двухуровневая система порождения вычислительных потоков},
  pdfauthor={development\_system}
}

% Цветовая схема для листингов
\definecolor{codebg}{RGB}{248,248,248}
\definecolor{codeframe}{RGB}{200,200,200}
\definecolor{keyword}{RGB}{0,0,180}
\definecolor{comment}{RGB}{100,120,100}
\definecolor{string}{RGB}{160,30,30}
\definecolor{builtin}{RGB}{130,0,130}

% Настройка listings
\lstdefinelanguage{Racket}{
  keywords={define, lambda, let, let*, letrec, cond, if, begin, and, or, not,
            when, unless, match, for, for/list, for/or, for/fold, map, filter,
            require, provide, module+, parameterize, with-handlers, values,
            string=?, string-prefix?, string-append, string-trim, string-join,
            hash, hash-ref, hash-has-key?, hash-set!, hash-remove!, make-hash,
            list, cons, car, cdr, first, second, third, fourth, fifth, null?,
            format, display, displayln, fprintf, newline, number->string,
            string->number, path->string, build-path, file-exists?, directory-exists?},
  morestring=[b]",
  morecomment=[l]{;;;},
  morecomment=[l]{;;},
  morecomment=[l]{;},
  sensitive=true
}

\lstdefinelanguage{CommonLisp}{
  keywords={defun, defpackage, defvar, defparameter, in-package, use-package,
            let, let*, labels, flet, cond, if, when, unless, and, or, not,
            progn, multiple-value-bind, values, restart-case, handler-bind,
            error, format, list, cons, car, cdr, first, second, string,
            string=, string/=, typep, declare, ignore},
  morestring=[b]",
  morecomment=[l]{;;;},
  morecomment=[l]{;;},
  morecomment=[l]{;},
  sensitive=nil
}

\lstset{
  backgroundcolor=\color{codebg},
  frame=single,
  rulecolor=\color{codeframe},
  basicstyle=\ttfamily\small,
  keywordstyle=\color{keyword}\bfseries,
  commentstyle=\color{comment}\itshape,
  stringstyle=\color{string},
  numbers=left,
  numberstyle=\tiny\color{gray},
  numbersep=8pt,
  breaklines=true,
  breakatwhitespace=true,
  tabsize=2,
  showstringspaces=false,
  captionpos=b,
  xleftmargin=1em,
  xrightmargin=0.5em
}

% Красивые заголовки секций
\titleformat{\section}{\large\bfseries}{\thesection.}{0.7em}{}[\vspace{-0.3em}\rule{\linewidth}{0.4pt}]
\titleformat{\subsection}{\normalsize\bfseries}{\thesubsection.}{0.5em}{}

% Колонтитулы
\pagestyle{fancy}
\fancyhf{}
\fancyhead[L]{\small\textit{development\_system}}
\fancyhead[R]{\small\textit{Двухуровневая система порождения вычислительных потоков}}
\fancyfoot[C]{\thepage}
\renewcommand{\headrulewidth}{0.3pt}

\setlength{\parindent}{0pt}
\setlength{\parskip}{0.7em}

% Окружение для ключевых понятий
\newmdenv[
  linewidth=0.5pt,
  linecolor=blue!40,
  backgroundcolor=blue!4,
  leftmargin=0,
  rightmargin=0,
  innerleftmargin=10pt,
  innerrightmargin=10pt,
  innertopmargin=8pt,
  innerbottommargin=8pt
]{conceptbox}

\begin{document}

% ============================================================
\begin{titlepage}
\vspace*{3cm}
\begin{center}
  {\LARGE\bfseries Двухуровневая система порождения\\[0.4em]
  и исполнения вычислительных потоков}\\[2em]
  {\large Архитектура, реализация и ключевые инженерные решения}\\[3em]
  {\normalsize\itshape
    Технический отчёт\\[0.5em]
    Март 2026}\\[4em]
  \begin{abstract}
    \normalsize
    В настоящей работе описана двухуровневая система \texttt{development\_system},
    предназначенная для интерактивного порождения и исполнения вычислительных потоков
    на основе больших языковых моделей. Система реализована в виде двух взаимодействующих
    Telegram-ботов: оркестратора~(\emph{Мастер}) на языке Racket и исполнителя
    (\emph{Подмастерье}), сочетающего Racket с Common~Lisp. Ключевой элемент системы —
    конвейер порождения: Мастер принимает от пользователя описание задачи на естественном
    языке, формирует структурированное задание, вызывает внешний агент Claude~Code,
    получает синтаксически корректный Common~Lisp, компилирует его средствами SBCL и
    сохраняет как именованный \emph{поток}. Подмастерье загружает активный поток и
    исполняет его для каждого пользовательского запроса. Описаны архитектура, протоколы
    взаимодействия компонентов, система условий и перезапуска в Common~Lisp,
    а также механизмы комбинирования потоков.
  \end{abstract}
\end{center}
\end{titlepage}

\tableofcontents
\newpage

% ============================================================
\section{Введение}

Взаимодействие человека с вычислительной машиной через язык естественных запросов
остаётся нерешённой инженерной задачей. Существующие подходы, как правило, сводятся
к одному из двух полюсов: либо к \emph{чёрному ящику} — диалоговой системе, не
способной к программной автоматизации, либо к \emph{явному программированию} —
требующему от пользователя глубокого знания API. \texttt{development\_system}
предлагает иной подход: \textbf{порождение кода на основе диалога} с немедленной
верификацией и исполнением результата.

Центральное понятие системы — \emph{вычислительный поток}: самодостаточная программа
на Common~Lisp, принимающая задачу в виде строки и возвращающая строковый ответ.
Потоки порождаются по описанию на естественном языке, проверяются компилятором и
выполняются в изолированном процессе SBCL. Такой подход позволяет:

\begin{itemize}[noitemsep]
  \item сохранять и повторно использовать задокументированную логику;
  \item комбинировать потоки последовательно, параллельно или условно;
  \item отделять роль «архитектора» (Мастер) от роли «исполнителя» (Подмастерье).
\end{itemize}

Система написана в духе SICP~\cite{sicp}: минимализм, рекурсия вместо итерации,
маленькие чистые функции, паттерн-матчинг. Имена прикладных понятий — на русском языке.

\subsection{Участники системы}

\begin{table}[H]
\centering
\begin{tabular}{lllp{6cm}}
\toprule
\textbf{Компонент} & \textbf{Язык} & \textbf{Роль} \\
\midrule
Мастер & Racket & Оркестратор: диалог, порождение потоков, управление \\
Подмастерье & Racket + SBCL & Исполнитель: приём задачи, запуск потока, ответ \\
Поток & Common Lisp & Вычислительная единица с точкой входа \texttt{выполнить} \\
Claude Code & Go (внешний) & Агент порождения кода, вызывается как подпроцесс \\
\bottomrule
\end{tabular}
\caption{Компоненты системы}
\end{table}

% ============================================================
\section{Архитектура системы}

\subsection{Обзор взаимодействия}

Система работает следующим образом. Мастер принимает сообщение пользователя через
Telegram Bot~API и классифицирует его: системная команда, кнопка меню или свободный
текст. Для команды \emph{«Породить поток»} запускается конвейер генерации. По
завершении пользователь активирует созданный поток командой \texttt{/запустить},
и с этого момента все запросы к Подмастерью исполняются через данный поток.

\begin{figure}[H]
\centering
\begin{mdframed}[linecolor=gray!60,backgroundcolor=gray!5]
\ttfamily\small
\begin{verbatim}
Пользователь
  │
  ├── [Telegram] ──► Мастер (Racket)
  │                    │
  │                    ├── Claude Code (subprocess)
  │                    │     └── порождает .lisp
  │                    ├── SBCL (subprocess, compile-file)
  │                    │     └── верифицирует код
  │                    └── файловая система ──► активный.scm
  │
  └── [Telegram] ──► Подмастерье (Racket)
                       │
                       └── SBCL (subprocess, load + выполнить)
                             └── OpenRouter API
\end{verbatim}
\end{mdframed}
\caption{Схема потоков данных}
\end{figure}

\subsection{Файловая среда как протокол}

Взаимодействие между Мастером и Подмастерьем осуществляется \emph{исключительно
через файловую систему}. Никакого IPC, никаких сокетов. Мастер записывает файл
\texttt{потоки/активный.scm}, Подмастерье его читает при получении каждого запроса.

Формат файла — S-выражение:

\begin{lstlisting}[language=Racket, caption={Файл активного потока \texttt{активный.scm}}]
(активный
  (имя    "уточнение")
  (файл   "/path/to/потоки/уточнение.lisp")
  (пакет  "поток-уточнение"))
\end{lstlisting}

Такое решение обеспечивает полную независимость компонентов и тривиальную отладку:
состояние системы всегда доступно для инспекции как обычный текстовый файл.

% ============================================================
\section{Конфигурация системы: концепция «Души»}

\subsection{Структура \texttt{душа.scm}}

Поведение Мастера определяется конфигурационным файлом, хранящимся в виде
S-выражения~— \emph{«Души»}. Файл содержит системное указание (system prompt),
предпочтения модели, долговременную память о пользователе и стилевые директивы:

\begin{lstlisting}[language=Racket, caption={Конфигурация системы \texttt{душа.scm}}]
(душа
  (указание  "Ты помощник для создания вычислительных потоков
              на Common Lisp.")
  (модель    "claude-sonnet-4-20250514")
  (теплота   0.3)
  (память    ((предпочитает . "рекурсию а не циклы")))
  (девиз     "краткость, однострочники, без англицизмов, SICP"))
\end{lstlisting}

\subsection{Операции над душой: \texttt{dusha.rkt}}

Модуль \texttt{dusha.rkt} реализует типичный функциональный CRUD для
неизменяемых S-выражений. Доступ к полям через паттерн-матчинг:

\begin{lstlisting}[language=Racket, caption={Извлечение поля из S-выражения (dusha.rkt)}]
(define (получить-поле душа имя)
  (match душа
    [(list 'душа поля ...)
     (for/or ([поле (in-list поля)])
       (match поле
         [(list (== имя) значение) значение]
         [_ #f]))]
    [_ #f]))
\end{lstlisting}

Обновление поля реализовано как структурное преобразование без мутации:

\begin{lstlisting}[language=Racket, caption={Функциональное обновление поля души}]
(define (заменить-поле душа имя значение)
  (match душа
    [(list 'душа поля ...)
     (cons 'душа
           (for/list ([поле (in-list поля)])
             (match поле
               [(list (== имя) _) (list имя значение)]
               [другое другое])))]))
\end{lstlisting}

Паттерн \texttt{(== имя)} обеспечивает структурное равенство при матчинге символа.

% ============================================================
\section{Модуль интеграции с языковой моделью: \texttt{llm.rkt}}

Для диалогового режима (свободный текст пользователя) Мастер обращается к
языковой модели через OpenRouter~API. Модуль \texttt{llm.rkt} инкапсулирует
этот вызов в одну функцию:

\begin{lstlisting}[language=Racket, caption={Ядро интеграции с LLM (llm.rkt)}]
(define (спросить-модель указание модель теплота сообщения)
  (let* ((м    (нормализовать-модель модель))
         (тело (jsexpr->bytes
                (hash 'model м
                      'temperature теплота
                      'messages
                      (cons (hash 'role "system" 'content указание)
                            сообщения))))
         (заголовки
          (list (string-append "Authorization: Bearer "
                               (ключ-openrouter))
                "Content-Type: application/json"))
         (порт (post-pure-port (string->url api-url) тело заголовки))
         (ответ (bytes->jsexpr (port->bytes порт))))
    (close-input-port порт)
    (let* ((варианты (hash-ref ответ 'choices #f))
           (первый   (and варианты
                          (> (vector-length варианты) 0)
                          (vector-ref варианты 0)))
           (сообщ    (and первый (hash-ref первый 'message #f)))
           (текст    (and сообщ  (hash-ref сообщ  'content #f))))
      (or текст "Нет ответа от модели."))))
\end{lstlisting}

Ключевые решения в реализации:

\begin{enumerate}[noitemsep]
  \item \textbf{Нормализация модели}: функция \texttt{нормализовать-модель}
    приводит краткие имена вида \texttt{"claude-sonnet-4"} к полным
    \texttt{"anthropic/claude-sonnet-4"}, ожидаемым OpenRouter.
  \item \textbf{JSON в Racket}: библиотека \texttt{racket/json} представляет
    JSON-массивы как \emph{списки}, а не векторы. Системный промпт
    добавляется через \texttt{cons}, формируя список сообщений без
    лишних аллокаций.
  \item \textbf{Идемпотентность}: функция не хранит состояния и не выполняет
    побочных эффектов кроме HTTP-запроса.
\end{enumerate}

% ============================================================
\section{Конвейер порождения потоков: \texttt{potoki.rkt}}

Это центральный модуль системы, реализующий полный жизненный цикл потока:
\textbf{описание → задание → код → верификация → хранение}.

\subsection{Вызов Claude Code как подпроцесса}

Для генерации кода используется инструмент Claude Code, запускаемый как
дочерний процесс через стандартный API подпроцессов Racket:

\begin{lstlisting}[language=Racket, caption={Вызов Claude Code (potoki.rkt)}]
(define (спросить-клод вопрос)
  (let ((путь-клод (find-executable-path "claude")))
    (unless путь-клод
      (error 'спросить-клод "Программа claude не найдена в PATH"))
    (let-values ([(процесс выход вход ошибки)
                  (subprocess #f #f #f
                              путь-клод "--print" вопрос)])
      (close-output-port вход)
      (let ((ответ (port->string выход)))
        (close-input-port выход)
        (close-input-port ошибки)
        (subprocess-wait процесс)
        (let ((код (subprocess-status процесс)))
          (if (zero? код) ответ
              (error 'спросить-клод
                     "Claude Code завершился с кодом ~a" код)))))))
\end{lstlisting}

Флаг \texttt{--print} переводит Claude Code в неинтерактивный режим: агент
получает запрос на stdin, выполняет его и завершается. Это делает вызов
прозрачным и детерминированным с точки зрения системы.

\subsection{Формирование задания}

Мастер формирует задание для Claude Code, объединяя контекст «Души» с
пользовательским описанием и явными техническими требованиями:

\begin{lstlisting}[language=Racket, caption={Формирование задания для генератора кода}]
(define (сформировать-задание душа описание)
  (let ((указание (or (получить-поле душа 'указание) ""))
        (девиз    (or (получить-поле душа 'девиз) "")))
    (string-append
     "Ты генератор вычислительных потоков на Common Lisp.\n"
     "Системное указание: " указание "\n"
     "Девиз: " девиз "\n\n"
     "Задача: " описание "\n\n"
     "Требования к коду:\n"
     "- В начале после defpackage: "
     "  (ql:quickload '(\"drakma\" \"cl-json\") :silent t)\n"
     "- HTTP: drakma:http-request; JSON: cl-json\n"
     "- Getenv: uiop:getenv (OPENROUTER_API_KEY)\n"
     "- Пакет :поток-<имя>; функция (defun выполнить (задача) ...)\n"
     "- Рекурсия вместо loop; функции до 15 строк; без CLOS\n"
     "- Верни ТОЛЬКО код Common Lisp\n")))
\end{lstlisting}

\subsection{Извлечение кода и определение имени пакета}

Ответ Claude Code может содержать markdown-обёртку (\texttt{```lisp...```}).
Функция \texttt{извлечь-код} реализует конечный автомат над строками ответа:

\begin{lstlisting}[language=Racket, caption={Извлечение кода из ответа (конечный автомат)}]
(define (извлечь-код ответ)
  (let ((строки (string-split ответ "\n")))
    (let ищи ((ост строки) (внутри #f) (собрано '()))
      (cond
        [(null? ост)
         (if (null? собрано)
             (string-trim ответ)        ; код без обёртки
             (string-join (reverse собрано) "\n"))]
        [(and (not внутри)
              (let ((с (string-trim (car ост))))
                (or (string-prefix? с "```lisp")
                    (string-prefix? с "```common-lisp"))))
         (ищи (cdr ост) #t собрано)]   ; вошли в блок
        [(and внутри
              (string-prefix? (string-trim (car ост)) "```"))
         (ищи (cdr ост) #f собрано)]   ; вышли из блока
        [внутри
         (ищи (cdr ост) #t (cons (car ост) собрано))]
        [else
         (ищи (cdr ост) #f собрано)]))))
\end{lstlisting}

\begin{conceptbox}
\textbf{Важное замечание о \texttt{string-prefix?} в Racket.}
Функция принимает аргументы в порядке \texttt{(string-prefix?~строка~префикс)},
то есть возвращает \texttt{\#t}, если второй аргумент является префиксом первого.
Это противоположно интуитивному порядку, принятому в большинстве языков, и
является источником типичной ошибки при портировании кода.
\end{conceptbox}

Имя пакета извлекается регулярным выражением Perl-совместимого формата (\texttt{\#px}):

\begin{lstlisting}[language=Racket, caption={Извлечение имени пакета}]
(define (имя-пакета код)
  (let ((совп (regexp-match #px":поток-([^\\s()]+)" код)))
    (if совп
        (string-trim (cadr совп))
        "безымянный")))
\end{lstlisting}

Использование \texttt{\#px} (PCRE) вместо \texttt{\#rx} (базовые регулярные выражения)
принципиально: только PCRE поддерживает мета-символы \texttt{\\s} и \texttt{\\w},
позволяющие корректно обработать пробельные символы и переносы строк.

\subsection{Верификация через компилятор SBCL}

Перед сохранением код компилируется в SBCL для проверки синтаксической корректности.
Подпроцесс запускается с предварительной загрузкой Quicklisp и зависимостей:

\begin{lstlisting}[language=Racket, caption={Верификация кода в SBCL (potoki.rkt)}]
(define (проверить-поток путь)
  (let* ((путь-sbcl  (find-executable-path "sbcl"))
         (ql-init    (path->string
                      (build-path (find-system-path 'home-dir)
                                  "quicklisp" "setup.lisp")))
         (загр-ql    (format "(when (probe-file ~s) (load ~s :verbose nil))"
                             ql-init ql-init))
         (загр-зав   (string-append
                      "(when (find-package :ql)"
                      "  (funcall (intern \"QUICKLOAD\" :ql)"
                      "   '(\"drakma\" \"cl-json\" \"uiop\")"
                      "   :silent t))"))
         (выражение  (format
                      "(multiple-value-bind (fasl warn fail)"
                      "    (compile-file ~s)"
                      "  (if fail (sb-ext:exit :code 1)"
                      "           (sb-ext:exit :code 0)))"
                      (path->string путь))))
    (let-values ([(проц выход вход ошибки)
                  (subprocess #f #f #f путь-sbcl
                              "--noinform" "--non-interactive"
                              "--eval" загр-ql
                              "--eval" загр-зав
                              "--eval" выражение)])
      (close-output-port вход)
      (let ((вывод (port->string выход))
            (ош    (port->string ошибки)))
        (close-input-port выход)
        (close-input-port ошибки)
        (subprocess-wait проц)
        (values (zero? (subprocess-status проц))
                (string-append вывод ош))))))
\end{lstlisting}

Функция \texttt{compile-file} возвращает три значения: путь к \texttt{.fasl}-файлу,
флаг предупреждений и флаг ошибок. Именно третье значение (\texttt{fail}) определяет
результат верификации. Код завершения подпроцесса транслируется обратно в булевое
значение.

% ============================================================
\section{Оркестратор: \texttt{мастер/main.rkt}}

\subsection{Reply-клавиатура}

Пользовательский интерфейс реализован через Telegram Reply~Keyboard — постоянную
клавиатуру, отображаемую в нижней части экрана мессенджера. Клавиша генерирует
текстовое сообщение, которое обрабатывается в обычном диспетчере:

\begin{lstlisting}[language=Racket, caption={Reply-клавиатура (main.rkt)}]
(define (кнопка-reply текст)
  (hash 'text текст))

(define (главное-меню-reply)
  (list
   (list (кнопка-reply "🔧 Породить поток")
         (кнопка-reply "📋 Мои потоки"))
   (list (кнопка-reply "ℹ️ Состояние")
         (кнопка-reply "⏹ Остановить"))
   (list (кнопка-reply "🧠 Душа")
         (кнопка-reply "💾 Память"))
   (list (кнопка-reply "🕐 Время"))))

(define (послать-с-кнопками ид-чата текст)
  (послать-запрос "sendMessage"
    (jsexpr->string
      (hash 'chat_id      ид-чата
            'text         текст
            'reply_markup
            (hash 'keyboard        (главное-меню-reply)
                  'resize_keyboard #t
                  'persistent      #t)))))
\end{lstlisting}

Флаг \texttt{persistent} обеспечивает постоянное присутствие клавиатуры даже
после перезапуска клиента.

\subsection{Диспетчер сообщений}

Диспетчер сопоставляет входящий текст с командами, включая русскоязычные кнопки меню:

\begin{lstlisting}[language=Racket, caption={Диспетчер команд}]
(define (сформировать-ответ текст)
  (cond
    [(or (string=? текст "/start")
         (string=? текст "/меню"))         #f]  ; показать меню
    [(or (string=? текст "🕐 Время")
         (string=? текст "/время"))        (формат-время)]
    [(or (string=? текст "ℹ️ Состояние")
         (string=? текст "/состояние"))    (обработать-состояние)]
    [(or (string=? текст "📋 Мои потоки")
         (string=? текст "/потоки"))       (обработать-потоки)]
    [(or (string=? текст "⏹ Остановить")
         (string=? текст "/остановить"))   (обработать-остановить)]
    [(string-prefix? текст "/породить ")
     (обработать-породить (substring текст 10))]
    [else (обработать-диалог текст)]))
\end{lstlisting}

\subsection{Управление состоянием «ожидания ввода»}

Кнопка «Породить поток» требует последующего ввода описания. Для этого используется
мутируемая хеш-таблица, отслеживающая чаты, находящиеся в режиме ввода:

\begin{lstlisting}[language=Racket, caption={Двухшаговый ввод через состояние}]
(define *ожидает-породить* (make-hash))

;; В цикле опроса:
(cond
  [(string=? текст "🔧 Породить поток")
   (hash-set! *ожидает-породить* ид-чата #t)
   (послать-сообщение ид-чата
     "Опишите задачу для нового потока:")]
  [(hash-ref *ожидает-породить* ид-чата #f)
   (hash-remove! *ожидает-породить* ид-чата)
   (let ((р (обработать-породить текст)))
     (послать-с-кнопками ид-чата р))]
  [else ...])
\end{lstlisting}

\subsection{Авторизация}

Оба бота ограничены одним администратором. Идентификатор задаётся через переменную
окружения \texttt{ADMIN\_CHAT\_ID}:

\begin{lstlisting}[language=Racket, caption={Проверка авторизации}]
(define (ид-админа)
  (let ((s (getenv "ADMIN_CHAT_ID")))
    (and s (string->number s))))

(define (админ? ид-чата)
  (let ((а (ид-админа)))
    (or (not а) (= ид-чата а))))

;; В цикле опроса: игнорируем чужих молча
(when (and (админ? ид-чата) (eq? тип 'message))
  ... ; обработка)
\end{lstlisting}

Если переменная не задана, авторизация отключена (\texttt{(not а)} = \texttt{\#t}).
Это позволяет запустить систему без конфигурации в режиме разработки.

% ============================================================
\section{Исполнитель: \texttt{подмастерье/исполнитель.lisp}}

Common Lisp исполнитель реализует загрузку потока, его верификацию,
выполнение задачи и обработку условий. Это единственный компонент системы
на Common Lisp (не считая самих потоков).

\subsection{Инициализация Quicklisp}

Поскольку SBCL запускается как чистый процесс без пользовательской инициализации,
необходима явная загрузка Quicklisp перед выполнением потока:

\begin{lstlisting}[language=CommonLisp, caption={Инициализация Quicklisp в исполнителе}]
(defun обеспечить-quicklisp ()
  "Загружает Quicklisp если ещё не загружен."
  (unless (find-package :ql)
    (let ((ql-init (merge-pathnames "quicklisp/setup.lisp"
                                    (user-homedir-pathname))))
      (when (probe-file ql-init)
        (load ql-init :verbose nil :print nil))))
  (when (find-package :ql)
    (funcall (intern "QUICKLOAD" :ql)
             '("drakma" "cl-json" "flexi-streams" "uiop")
             :silent t)))
\end{lstlisting}

Использование \texttt{intern} вместо прямого символа \texttt{ql:quickload} обязательно:
пакет \texttt{:ql} ещё не существует в момент компиляции файла.

\subsection{Загрузка потока}

\begin{lstlisting}[language=CommonLisp, caption={Загрузка потока с обработкой условий}]
(defun загрузить-поток (путь)
  "Загружает файл потока в текущий образ."
  (restart-case
      (progn
        (unless (probe-file путь)
          (error "Файл потока не найден: ~a" путь))
        (обеспечить-quicklisp)
        (load путь :verbose nil :print nil)
        t)
    (пропустить ()
      :report "Пропустить загрузку потока"
      nil)))
\end{lstlisting}

\subsection{Поиск точки входа}

Потоки могут иметь произвольные имена пакетов, поэтому точка входа ищется
по соглашению об именовании — функция \texttt{ВЫПОЛНИТЬ} в заданном пакете:

\begin{lstlisting}[language=CommonLisp, caption={Динамический поиск точки входа}]
(defun найти-точку-входа (имя-пакета)
  "Ищет функцию ВЫПОЛНИТЬ в указанном пакете."
  (let ((пакет (find-package (string-upcase имя-пакета))))
    (when пакет
      (let ((символ (find-symbol "ВЫПОЛНИТЬ" пакет)))
        (when (and символ (fboundp символ))
          (symbol-function символ))))))
\end{lstlisting}

\texttt{string-upcase} необходим: в Common Lisp символы хранятся в верхнем регистре,
а имена пакетов в метафайлах написаны строчными буквами.

% ============================================================
\section{Анатомия вычислительного потока}

Потоки генерируются по единому шаблону и имеют предсказуемую структуру.
Рассмотрим два реальных потока, порождённых системой.

\subsection{Поток «уточнение»: итеративное решение задач}

Поток реализует паттерн \emph{iterative refinement}: каждый шаг улучшает ответ
предыдущего, останавливаясь при достижении критерия качества или по исчерпании
лимита итераций.

\begin{lstlisting}[language=CommonLisp, caption={Поток уточнения: ядро алгоритма}]
(defun уточнить (задача предыдущий шаг)
  "Рекурсивно уточняет ответ. Шаг от 1 до 5."
  (if (> шаг 5)
      предыдущий
      (let* ((подсказка
               (if (string= предыдущий "")
                   (format nil "Задача: ~a" задача)
                   (format nil
                     "Задача: ~a~%Предыдущий ответ:~%~a~%~%Уточни."
                     задача предыдущий)))
             (история (list (наставление)
                            (сообщение "user" подсказка)))
             (ответ (послать-запрос история)))
        (if (достаточно-ли ответ)
            ответ
            (уточнить задача ответ (1+ шаг))))))
\end{lstlisting}

Критерий остановки:

\begin{lstlisting}[language=CommonLisp, caption={Критерий качества ответа}]
(defun достаточно-ли (ответ)
  "Проверяет признаки завершённости в тексте."
  (and (> (length ответ) 20)
       (or (search "ГОТОВО" ответ)
           (search "Итог:" ответ)
           (search "Ответ:" ответ))))
\end{lstlisting}

Точка входа оборачивает рекурсию в \texttt{restart-case}:

\begin{lstlisting}[language=CommonLisp, caption={Точка входа потока уточнения}]
(defun выполнить (задача)
  "Итеративно решает задачу за 1–5 шагов."
  (restart-case
      (уточнить задача "" 1)
    (вернуть-ошибку (e)
      :report "Вернуть описание ошибки"
      (format nil "Ошибка: ~a" e))))
\end{lstlisting}

\subsection{Поток «исполнитель»: порождение и выполнение кода}

Поток реализует полный цикл: получить задачу → попросить модель написать Python-код
→ очистить его от markdown-обёртки → запустить через \texttt{uiop:run-program}.

\begin{lstlisting}[language=CommonLisp, caption={Исполнение Python-кода через uiop}]
(defun исполнить-python (код)
  (restart-case
      (multiple-value-bind (вывод код-ошибки статус)
          (uiop:run-program
           (list "python3" "-c" код)
           :output        :string
           :error-output  :string
           :ignore-error-status t)
        (declare (ignore статус))
        (if (and код-ошибки (string/= код-ошибки ""))
            (values nil код-ошибки)
            (values вывод nil)))
    (вернуть-ошибку (e)
      (values nil (format nil "Ошибка: ~A" e)))))
\end{lstlisting}

Очистка кода от markdown удаляет строки, начинающиеся с \texttt{```}:

\begin{lstlisting}[language=CommonLisp, caption={Рекурсивная очистка markdown-обёртки}]
(defun убрать-обёртку (строки)
  (cond ((null строки) nil)
        ((начинается-с "```" (car строки))
         (убрать-обёртку (cdr строки)))
        (t (cons (car строки)
                 (убрать-обёртку (cdr строки))))))

(defun начинается-с (префикс строка)
  (and (>= (length строка) (length префикс))
       (string= префикс строка :end2 (length префикс))))
\end{lstlisting}

Точка входа объединяет все шаги:

\begin{lstlisting}[language=CommonLisp, caption={Точка входа потока исполнителя}]
(defun выполнить (задача)
  (restart-case
      (let* ((код-python (запросить-модель задача))
             (чистый-код (очистить-код код-python)))
        (multiple-value-bind (результат ошибка)
            (исполнить-python чистый-код)
          (or ошибка результат)))
    (отказ (причина)
      (format nil "Отказ: ~A" причина))))
\end{lstlisting}

% ============================================================
\section{Комбинирование потоков: \texttt{kombinirovanie.rkt}}

Мастер поддерживает три базовые схемы комбинирования потоков и произвольную
четвёртую, генерируемую по описанию.

\subsection{Последовательная цепочка}

Порождает новый поток-обёртку, последовательно передающий результат:

\begin{lstlisting}[language=Racket, caption={Генерация кода цепочки потоков}]
(define (шаблон-цепочки имена пакеты файлы)
  (string-append
   "(defpackage :поток-" (string-join имена "-") " (:use :cl))\n"
   "(in-package :поток-" (string-join имена "-") ")\n\n"
   ;; загрузка всех потоков
   (apply string-append
          (map (lambda (ф)
                 (format "(load ~s :verbose nil :print nil)\n" ф))
               файлы))
   "\n(defun выполнить (задача)\n"
   "  (let цикл ((вх задача)\n"
   "             (ост (list " (string-join
                               (map (lambda (п)
                                      (string-append "'" п ":выполнить"))
                                    пакеты)
                               " ") ")))\n"
   "    (if (null ост) вх\n"
   "        (цикл (funcall (car ост) вх) (cdr ост)))))\n"))
\end{lstlisting}

% ============================================================
\section{Протокол взаимодействия с Telegram Bot API}

\subsection{Long polling}

Оба бота используют long polling: запрос \texttt{getUpdates} с \texttt{timeout=30}
блокируется на сервере Telegram до появления нового обновления или истечения тайм-аута.

\begin{lstlisting}[language=Racket, caption={Long polling в Racket}]
(define (получить-обновления сдвиг)
  (let* ((адрес (string->url
                 (string-append (адрес-api "getUpdates")
                                "?timeout=30&offset="
                                (number->string сдвиг))))
         (порт (get-pure-port адрес)))
    (begin0 (port->string порт)
            (close-input-port порт))))
\end{lstlisting}

Параметр \texttt{offset} содержит \texttt{update\_id + 1} последнего обработанного
обновления, что исключает повторную обработку.

\subsection{Разбор обновлений}

Парсер обрабатывает два типа событий: обычные сообщения и ответы на callback-кнопки.
Результат унифицирован в единый формат события:

\begin{lstlisting}[language=Racket, caption={Парсер обновлений: унификация событий}]
(define (извлечь-события обновления)
  (filter-map
   (lambda (обн)
     (cond
       [(hash-has-key? обн 'message)
        (let* ((сообщ   (hash-ref обн 'message))
               (ид-обн  (hash-ref обн 'update_id 0))
               (чат     (hash-ref сообщ 'chat (hash)))
               (ид-чата (hash-ref чат 'id 0))
               (текст   (hash-ref сообщ 'text "")))
          (list ид-обн ид-чата 'message текст))]
       [(hash-has-key? обн 'callback_query)
        (let* ((cb      (hash-ref обн 'callback_query))
               (ид-обн  (hash-ref обн 'update_id 0))
               (чат     (hash-ref
                          (hash-ref cb 'message (hash))
                          'chat (hash)))
               (ид-чата (hash-ref чат 'id 0))
               (cb-id   (hash-ref cb 'id ""))
               (данные  (hash-ref cb 'data "")))
          (list ид-обн ид-чата 'callback cb-id данные))]
       [else #f]))
   обновления))
\end{lstlisting}

% ============================================================
\section{Инженерные уроки и ключевые ошибки}

В процессе разработки были выявлены несколько нетривиальных особенностей
применяемых инструментов:

\subsection{Порядок аргументов \texttt{string-prefix?} в Racket}

\begin{conceptbox}
Функция \texttt{(string-prefix? s1 s2)} в \texttt{racket/string}
возвращает \texttt{\#t} если \emph{s2} является префиксом \emph{s1}.
Иными словами, первый аргумент — строка, второй — искомый префикс.
Это противоположно семантике аналогичных функций в Python, Haskell и Go.
\end{conceptbox}

Проверка: \texttt{(string-prefix? "/породить тест" "/породить ")} → \texttt{\#t}.

\subsection{JSON-массивы в Racket: списки, не векторы}

Библиотека \texttt{racket/json} представляет JSON-массивы как списки Racket,
а не векторы. Передача вектора в \texttt{jsexpr->bytes} вызывает ошибку типа.
Корректное формирование массива сообщений:

\begin{lstlisting}[language=Racket, caption={Правильное формирование JSON-массива}]
;; Неправильно (vector):
'messages (list->vector ...)   ; → ошибка jsexpr->bytes

;; Правильно (list):
'messages (cons (hash 'role "system" 'content указание)
                сообщения)
\end{lstlisting}

\subsection{Регулярные выражения: \texttt{\#rx} vs \texttt{\#px}}

В Racket существуют два типа регулярных выражений:
\texttt{\#rx} (POSIX-совместимые) не обрабатывают \texttt{\\s}, \texttt{\\w} и
аналогичные мета-символы. Для работы с ними необходим \texttt{\#px} (PCRE):

\begin{lstlisting}[language=Racket]
;; #rx не понимает \s:
(regexp-match #rx":поток-([^\s()]+)" код)   ; ошибка

;; #px работает корректно:
(regexp-match #px":поток-([^\\s()]+)" код)  ; OK
\end{lstlisting}

\subsection{Динамическая загрузка Quicklisp в Common Lisp}

При вызове \texttt{ql:quickload} в файле, загружаемом до инициализации Quicklisp,
компилятор сигнализирует об ошибке «Package QL does not exist». Обходное решение —
\texttt{intern}:

\begin{lstlisting}[language=CommonLisp, caption={Безопасная загрузка через intern}]
;; Неправильно (ошибка компиляции до загрузки Quicklisp):
(ql:quickload '("drakma") :silent t)

;; Правильно (откладывает разрешение символа на время выполнения):
(funcall (intern "QUICKLOAD" :ql) '("drakma") :silent t)
\end{lstlisting}

% ============================================================
\section{Переменные окружения и конфигурация}

\begin{table}[H]
\centering
\begin{tabular}{lll}
\toprule
\textbf{Переменная} & \textbf{Компонент} & \textbf{Назначение} \\
\midrule
\texttt{MASTER\_BOT\_TOKEN}     & Мастер      & Токен Telegram-бота Мастера \\
\texttt{APPRENTICE\_BOT\_TOKEN} & Подмастерье & Токен Telegram-бота Подмастерья \\
\texttt{OPENROUTER\_API\_KEY}   & Оба         & Ключ API OpenRouter \\
\texttt{ADMIN\_CHAT\_ID}        & Оба         & Telegram ID единственного пользователя \\
\bottomrule
\end{tabular}
\caption{Переменные окружения системы}
\end{table}

% ============================================================
\section{Метафайлы потоков: формат \texttt{.мета}}

Каждый поток сопровождается файлом метаданных в формате S-выражения.
Это решение продолжает сквозную тему проекта: состояние системы — всегда
читаемый текст.

\begin{lstlisting}[language=Racket, caption={Формат файла метаданных потока}]
(мета
  (имя         "уточнение")
  (создан      "2026-03-13T10:17:12")
  (описание    "Итеративное решение задач через языковую модель")
  (зависимости ())
  (входит-в    ())   ; список потоков-агрегаторов
  (состояние   :готов))  ; :готов | :работает | :ошибка
\end{lstlisting}

Доступ к полям метафайла через паттерн-матчинг аналогичен работе с «Душой»:

\begin{lstlisting}[language=Racket, caption={Чтение метаданных потока}]
(define (поле-меты мета имя)
  (match мета
    [(list 'мета поля ...)
     (for/or ([поле (in-list поля)])
       (match поле
         [(list (== имя) значение) значение]
         [_ #f]))]
    [_ #f]))

;; Пример использования:
(define (описание-потока имя)
  (let* ((путь-мета (build-path каталог
                                (string-append имя ".мета")))
         (мета      (прочитать-мету путь-мета)))
    (or (and мета (поле-меты мета 'описание))
        "без описания")))
\end{lstlisting}

% ============================================================
\section{Цикл опроса: обработка ошибок и перезапуск}

\subsection{Устойчивость к сетевым сбоям}

Главный цикл опроса в обоих ботах оборачивает каждую итерацию в
\texttt{with-handlers}. При любом исключении — сетевом тайм-ауте,
ошибке парсинга JSON, недоступности API — цикл засыпает на пять секунд
и продолжает работу с тем же сдвигом:

\begin{lstlisting}[language=Racket, caption={Устойчивый цикл опроса с перезапуском}]
(define (запустить-опрос)
  (загрузить-душу)
  (displayln "Мастер запущен. Ожидаю сообщения...")
  (let цикл ((сдвиг 0))
    (with-handlers
        ([exn:fail?
          (lambda (e)
            (displayln
              (format "Ошибка: ~a" (exn-message e)))
            (sleep 5)
            (цикл сдвиг))])   ; повторяем с тем же сдвигом
      (let* ((ответ    (получить-обновления сдвиг))
             (события  (извлечь-события
                          (разобрать-обновления ответ))))
        (let обр ((ост события) (макс сдвиг))
          (if (null? ост)
              (цикл макс)
              (let* ((событие (car ост))
                     (ид-обн  (first событие))
                     ...)
                ...
                (обр (cdr ост)
                     (max макс (+ ид-обн 1))))))))))
\end{lstlisting}

Ключевой инвариант: \texttt{сдвиг} не обновляется до успешной обработки
всех сообщений текущего батча. Это гарантирует, что при падении в середине
батча все сообщения будут обработаны повторно.

\subsection{Защита от чужих пользователей}

Проверка авторизации происходит до любой обработки события, чтобы не
расходовать ресурсы LLM на нежелательные запросы:

\begin{lstlisting}[language=Racket, caption={Раннее отклонение неавторизованных}]
(when (and (админ? ид-чата) (eq? тип 'message))
  ;; только здесь происходит реальная обработка
  (let ((текст (fourth событие)))
    ...))
;; чужие молча игнорируются — никакого ответа
\end{lstlisting}

% ============================================================
\section{Запуск потока в Подмастерье: полный путь данных}

Цепочка вызовов при исполнении пользовательского запроса в Подмастерье
охватывает несколько уровней абстракции: от Telegram до SBCL-подпроцесса.

\begin{lstlisting}[language=Racket, caption={Исполнение потока: Racket-часть Подмастерья}]
(define (выполнить-в-sbcl путь-потока имя-пакета задача)
  (let* ((путь-sbcl (find-executable-path "sbcl"))
         (исп       (->string (путь-исполнителя)))
         (загрузка
          (format "(load ~s :verbose nil :print nil)" исп))
         (вызов
          (format
           "(let ((ф (funcall (intern \"ЗАПУСК\" \"ИСПОЛНИТЕЛЬ\")
                               ~s ~s ~s)))
              (princ ф) (finish-output)
              (sb-ext:exit :code 0))"
           (->string путь-потока) имя-пакета задача)))
    (let-values ([(проц выход вход ошибки)
                  (subprocess #f #f #f
                              путь-sbcl
                              "--noinform" "--non-interactive"
                              "--eval" загрузка
                              "--eval" вызов)])
      (close-output-port вход)
      (let ((вывод (port->string выход))
            (ош    (port->string ошибки)))
        (close-input-port выход)
        (close-input-port ошибки)
        (subprocess-wait проц)
        (if (zero? (subprocess-status проц))
            (string-trim вывод)
            (format "Ошибка SBCL: ~a" (string-trim ош)))))))
\end{lstlisting}

Функция \texttt{ЗАПУСК} в пакете \texttt{ИСПОЛНИТЕЛЬ} загружает поток,
находит точку входа и вызывает её:

\begin{lstlisting}[language=CommonLisp, caption={Функция ЗАПУСК в исполнителе (CL)}]
(defun запуск (путь-потока имя-пакета задача)
  "Загружает поток, находит точку входа, исполняет."
  (загрузить-поток путь-потока)
  (let ((точка-входа (найти-точку-входа имя-пакета)))
    (unless точка-входа
      (error "Функция ВЫПОЛНИТЬ не найдена в пакете ~a"
             имя-пакета))
    (выполнить-с-повтором точка-входа задача)))

(defun выполнить-с-повтором (функция задача)
  "Вызывает функцию потока с защитой от условий."
  (restart-case
      (funcall функция задача)
    (вернуть-задачу ()
      :report "Вернуть исходную задачу при ошибке"
      задача)))
\end{lstlisting}

% ============================================================
\section{Производительность и задержки}

\subsection{Источники задержки}

Время ответа системы складывается из нескольких компонентов:

\begin{table}[H]
\centering
\begin{tabular}{lrp{5.5cm}}
\toprule
\textbf{Этап} & \textbf{Типичное время} & \textbf{Примечание} \\
\midrule
Получение обновления (long poll) & $<$100 мс & при наличии сообщения \\
Запуск SBCL подпроцесса           & 1--3 с    & холодный старт \\
Загрузка Quicklisp + зависимостей & 2--5 с    & drakma, cl-json \\
Обращение к OpenRouter API        & 1--5 с    & зависит от модели \\
Итого (исполнение потока)         & 4--13 с   & \\
\midrule
Порождение потока (Claude Code)   & 15--60 с  & однократно \\
\bottomrule
\end{tabular}
\caption{Типичные задержки компонентов}
\end{table}

Наибольший вклад вносит запуск SBCL: каждый запрос к Подмастерью
создаёт новый процесс. Потенциальное улучшение — постоянный SBCL-образ
с предзагруженными зависимостями (\texttt{save-lisp-and-die}).

\subsection{Возможности оптимизации}

\begin{enumerate}
  \item \textbf{Образы SBCL}: сохранение образа с загруженными потоком
    и зависимостями сокращает холодный старт с 4--8 с до $<$100 мс.

  \item \textbf{Постоянный процесс-исполнитель}: вместо запуска нового
    SBCL для каждого запроса можно поддерживать долгоживущий процесс
    с SWANK-сервером для подачи задач.

  \item \textbf{Кеширование зависимостей}: предкомпиляция \texttt{.fasl}
    файлов устраняет повторную компиляцию \texttt{drakma} при каждом запросе.
\end{enumerate}

% ============================================================
\section{Заключение}

В работе описана система \texttt{development\_system}, реализующая концепцию
\emph{порождаемого кода как конфигурации}. Основные результаты:

\begin{enumerate}
  \item \textbf{Конвейер генерации}: описание на естественном языке
    преобразуется в верифицированный Common Lisp код за одно обращение
    к языковой модели.

  \item \textbf{Двухуровневая архитектура}: разделение ролей оркестратора
    (Racket, диалог и управление) и исполнителя (SBCL, вычисления) позволяет
    независимо масштабировать каждую компоненту.

  \item \textbf{Файловая среда как протокол}: отсутствие явного IPC между
    компонентами упрощает отладку и обеспечивает персистентность состояния.

  \item \textbf{Система условий Common Lisp}: паттерн \texttt{restart-case}
    обеспечивает структурированную обработку ошибок без исключений.

  \item \textbf{Стиль SICP}: рекурсия, маленькие функции, паттерн-матчинг
    и идиоматический Racket обеспечивают высокую читаемость при минимальном
    объёме кода.
\end{enumerate}

Открытые вопросы: персистентность истории диалога между сессиями,
параллельное исполнение потоков с разделением результатов, формальная
верификация потоков средствами ACL2.

% ============================================================
\section*{Приложение: структура файлов проекта}
\addcontentsline{toc}{section}{Приложение: структура файлов проекта}

\begin{lstlisting}[basicstyle=\ttfamily\footnotesize, numbers=none,
                   caption={Полная структура каталогов development\_system}]
development_system/
├── мастер/
│   ├── main.rkt           # Оркестратор: цикл опроса, меню, диспетчер
│   ├── dusha.rkt          # Конфигурация: чтение/запись S-выражений
│   ├── potoki.rkt         # Конвейер порождения и проверки потоков
│   ├── kombinirovanie.rkt # Комбинирование: цепочка, параллель, ветвление
│   ├── llm.rkt            # Интеграция с OpenRouter API
│   ├── душа.scm           # Конфигурационный файл системы
│   ├── история/           # Журналы диалогов (по дням, .txt)
│   └── потоки/
│       ├── активный.scm   # Активный поток для Подмастерья
│       ├── уточнение.lisp # Поток: итеративное уточнение ответа
│       ├── уточнение.мета # Метаданные потока
│       ├── исполнитель.lisp # Поток: генерация и выполнение кода
│       └── исполнитель.мета
│
├── подмастерье/
│   ├── main.rkt           # Бот-исполнитель: цикл опроса, диспетчер
│   └── исполнитель.lisp   # CL-ядро: загрузка, запуск, condition/restart
│
└── docs/
    ├── plans/
    │   └── TODO.md        # Исходное техническое задание
    └── paper.tex          # Настоящий документ
\end{lstlisting}

Общий объём кода системы (без генерируемых потоков):

\begin{table}[H]
\centering
\begin{tabular}{lrr}
\toprule
\textbf{Файл} & \textbf{Строк} & \textbf{Язык} \\
\midrule
\texttt{мастер/main.rkt}           & $\sim$330 & Racket \\
\texttt{мастер/potoki.rkt}         & $\sim$270 & Racket \\
\texttt{мастер/dusha.rkt}          & $\sim$120 & Racket \\
\texttt{мастер/kombinirovanie.rkt} & $\sim$160 & Racket \\
\texttt{мастер/llm.rkt}            & $\sim$40  & Racket \\
\texttt{подмастерье/main.rkt}      & $\sim$260 & Racket \\
\texttt{подмастерье/исполнитель.lisp} & $\sim$100 & Common Lisp \\
\midrule
\textbf{Итого}                     & $\sim$1280 & \\
\bottomrule
\end{tabular}
\caption{Объём кода системы}
\end{table}

\begin{thebibliography}{9}
\bibitem{sicp}
  Abelson, H., Sussman, G.J.
  \emph{Structure and Interpretation of Computer Programs}.
  MIT Press, 1996.

\bibitem{sbcl}
  SBCL Team.
  \emph{SBCL User Manual}.
  \url{https://www.sbcl.org/manual/}

\bibitem{racket}
  Flatt, M., PLT.
  \emph{Reference: Racket}.
  \url{https://docs.racket-lang.org/reference/}

\bibitem{openrouter}
  OpenRouter.
  \emph{OpenRouter API Documentation}.
  \url{https://openrouter.ai/docs}

\bibitem{telegram}
  Telegram.
  \emph{Bot API Reference}.
  \url{https://core.telegram.org/bots/api}
\end{thebibliography}

\end{document}
